\subsection{Actual harmonic restriction and polynomial units}\label{subsec:fc-line-harmonic-restriction}

\begin{lemma}[A nonvanishing linear logarithm is actually harmonic]\label{lem:fc-linear-unit-log-harmonic}
\lean{EnvelopeProperty.Berkovich.RelativeSpectrum.FCStandardAffinoid.isHarmonic_logLinear_of_nonvanishing}
\leanok
\uses{def:fc-actual-harmonic, lem:fc-inflation-linear}
If $T-a$ does not vanish at any actual point of a standard affinoid $U$,
its real logarithm is an actual harmonic function on $U$.
\end{lemma}

\begin{proof}
The rigid point $a$ is outside $U$. If it is outside the outer closed
disc, the ultrametric inequality makes $|T-a|$ constant throughout $U$.
Otherwise it lies in one of the actual open holes, say centred at $a_j$.
The strict inequality $|a-a_j|<r_j\le|T-a_j|_x$ gives
$|T-a|_x=|T-a_j|_x$ at every $x\in U$. Its logarithm is therefore one
of the literal original harmonic coordinates.
\end{proof}

\begin{lemma}[Actual harmonic restriction]\label{lem:fc-actual-harmonic-restriction}
\lean{EnvelopeProperty.Berkovich.RelativeSpectrum.FCStandardAffinoid.IsHarmonic.restrict, EnvelopeProperty.Berkovich.RelativeSpectrum.FCStandardAffinoid.isLineHarmonicOn_harmonicExpression}
\leanok
\uses{lem:fc-linear-unit-log-harmonic, lem:fc-harmonic-vector-space}
An actual harmonic function on a standard affinoid restricts to an actual
harmonic function on every actual standard subaffinoid. In particular,
an ambient harmonic expression is harmonic on the whole carrier in the
local test-affinoid sense.
\end{lemma}

\begin{proof}
Every original hole coordinate is nonvanishing on the original carrier,
hence on the subaffinoid. Apply the preceding lemma to each coordinate
and take its literal finite real linear combination and constant term.
\end{proof}

\begin{lemma}[Actual polynomial-unit logarithms]\label{lem:fc-polynomial-unit-log-harmonic}
\lean{EnvelopeProperty.Berkovich.RelativeSpectrum.isLineHarmonicOn_polynomialRealLog, EnvelopeProperty.Berkovich.RelativeSpectrum.isLineHarmonicOn_rationalRealLog}
\leanok
\uses{lem:fc-linear-unit-log-harmonic, lem:fc-harmonic-vector-space,
  def:fc-actual-line-harmonic-on}
The real logarithm of a polynomial is harmonic on every actual domain on
which its actual evaluations are all nonzero. The same holds for the real
logarithm of an actual polynomial ratio on its common unit locus.
\end{lemma}

\begin{proof}
Split the original polynomial over the algebraically closed field.
Nonvanishing of the whole polynomial implies nonvanishing of each linear
factor. Its real logarithm is the logarithm of the original leading
coefficient plus the sum of the original root logarithms with their actual
multiplicities. Apply the linear-unit lemma and actual harmonic vector
operations. A polynomial ratio is the actual difference of its two
polynomial logarithms on the common unit locus.
\end{proof}
